\setlength{\parindent}{0pt}

\section{Teoria dos Grafos}

\subsection{Teoremas}

\textbf{Euler:} Circuito Euleriano $\Leftrightarrow$ todos os vértices têm grau par. Caminho Euleriano $\Leftrightarrow$ exatamente 2 vértices de grau ímpar.

\textbf{Fórmula de Euler (planar):} $V-E+F=2$. Grafo planar simples: $E\leq3V{-}6$; planar bipartido: $E\leq2V{-}4$. Todo planar tem algum vértice de grau $\leq5$.

\textbf{Cayley:} Número de árvores geradoras rotuladas de $K_n$ é $n^{n-2}$.

\textbf{Kirchhoff:} Número de árvores geradoras = qualquer cofator da matriz Laplaciana $L=D-A$.

\textbf{König:} Em grafo bipartido: cobertura mínima de vértices $=$ emparelhamento máximo; conjunto independente máximo $= V -$ emparelhamento máximo.

\textbf{Hall:} Existe emparelhamento perfeito de $U$ $\Leftrightarrow$ $\forall S\subseteq U: |N(S)|\geq|S|$.

\textbf{Menger:} Número máximo de caminhos aresta-disjuntos entre $s$ e $t$ $=$ tamanho do corte mínimo $s$-$t$. Vale também para caminhos vértice-disjuntos.

\textbf{Dilworth:} Em poset, partição mínima em cadeias $=$ anti-cadeia máxima. (Reduz a emparelhamento no fecho transitivo.)

\textbf{Mirsky:} Em poset, partição mínima em anti-cadeias $=$ comprimento da maior cadeia.

\textbf{Prüfer:} Bijeção entre árvores rotuladas de $n$ vértices e sequências de comprimento $n{-}2$ sobre $\{1,\ldots,n\}$.

\textbf{Erdős-Gallai:} Sequência $d_1\geq\cdots\geq d_n$ é sequência de graus de algum grafo simples $\Leftrightarrow$ $\sum d_i$ é par e para todo $k$:
\[\sum_{i=1}^k d_i\leq k(k{-}1)+\sum_{i=k+1}^n\min(d_i,k)\]

\textbf{$k$-tough:} Grafo é $k$-tough se $\forall S\subseteq V, S\neq\emptyset$: $k\cdot c(G-S)\leq|S|$, onde $c(G-S)$ é o número de componentes conexas de $G-S$.

\subsection{Tipos de Grafo}

\noindent\begin{minipage}[t]{0.49\linewidth}
\includegraphics[width=\linewidth]{../../Macacário_Maratona_de_Programação/grafos_1.png}
\end{minipage}\hfill\begin{minipage}[t]{0.49\linewidth}
\includegraphics[width=\linewidth]{../../Macacário_Maratona_de_Programação/grafos_2.png}
\end{minipage}
